\section{Unique Crossing Times and Temporal Order}
\label{sec:crossing-order}

A two-wall transition gives one affine threshold hit on the positive support coordinate and one on the negative coordinate. To turn these hits into a temporal classification, the crossing parameter must be intrinsic rather than an arbitrary existential witness.

\subsection{Uniqueness}

Suppose \(x\neq y\). The affine map
\[
  \gamma_{x,y}(t)=x+t(y-x)
\]
is injective in its parameter. Lean proves this directly over the rationals by cancelling the nonzero factor \(y-x\).

\begin{theorem}[Unique affine threshold time]
Let \(x\neq y\). If
\[
  \gamma_{x,y}(t)=c
  \qquad\text{and}\qquad
  \gamma_{x,y}(s)=c,
\]
then
\[
  t=s.
\]
Consequently every straddling affine segment has a unique threshold crossing parameter in \([0,1]\cap\mathbb{Q}\).
\statusverified
\end{theorem}

For a two-wall belief transition write
\[
  t_+=\text{the unique positive-support crossing time},
\]
\[
  t_-=\text{the unique negative-support crossing time}.
\]
Both lie in the rational unit interval.

\subsection{Crossing-order trichotomy}

The linear order on \(\mathbb{Q}\) now yields an exhaustive classification:
\[
  t_+<t_-,
  \qquad
  t_+=t_-,
  \qquad
  t_-<t_+.
\]

\begin{theorem}[Affine crossing-order trichotomy]
Every two-wall conditionalized belief transition admits intrinsic positive and negative crossing times satisfying exactly one of
\[
  t_+<t_-,\qquad
  t_+=t_-,\qquad
  t_-<t_+.
\]
\statusverified
\end{theorem}

The three cases have immediate event readings:
\[
\begin{array}{ccl}
 t_+<t_- &:& \text{positive threshold wall first},\\
 t_+=t_- &:& \text{simultaneous crossing},\\
 t_-<t_+ &:& \text{negative threshold wall first}.
\end{array}
\]

\subsection{The simultaneous case}

The equality case has a geometric characterization stronger than mere coincidence of two stored parameters.

\begin{theorem}[Simultaneity as the wall intersection]
For a two-coordinate affine crossing pair,
\[
  t_+=t_-
\]
if and only if there exists one rational parameter \(t\) such that
\[
  \Ppos_t=c
  \qquad\text{and}\qquad
  \Pneg_t=c.
\]
Equivalently, the support path passes through the wall-intersection point
\[
  (c,c).
\]
\statusverified
\end{theorem}

Thus simultaneous diagonal motion is geometrically exceptional. If \(t_+\neq t_-\), there is a nonempty rational interval between the two wall events. The next section proves that this interval carries a forced categorical phase.

\paragraph{Why order matters.}
The endpoint projection forgets which threshold coordinate changes first. Yet for a diagonal transition both bits must change. The temporal order of those changes therefore contains information absent from the pair \((\text{source},\text{target})\). Crossing-order geometry recovers exactly that lost information.
